Nguồn: \href{https://www.facebook.com/truongquang.de/posts/1795899743799506}{Facebook tác giả}

Cách đây ít lâu chúng ta mất Văn Như Cương, một nhà giáo thiết tha với sự nghiệp đổi mới giáo dục, một tài năng sư phạm; nay chúng ta lại mất Phan Đình Diệu, nhà toán học tiên phong chuẩn bị cho nước ta tiến vào cách mạng 4.0, nhà tư tưởng có tầm nhìn cao rộng luôn gắn bó với thân phận dân tộc.

Phan Đình Diệu là một trí thức theo đúng nghĩa của từ này, tức là hoạt động trí tuệ vì đam mê khoa học, vì ham muốn cống hiến xây dựng xã hội chứ không vì lợi ích hay danh vọng cá nhân. Đó là một trí thức nghĩ sao nói vậy, không sợ cường quyền, không hùa theo kẻ mạnh. Ông là người không hề thay đổi sắc màu tư tưởng từ tuổi trẻ đến tuổi già. Điều đó một phần do môi trường sống qui định. Thân phụ ông từng tham gia các cuộc nổi dậy chống thực dân cai trị và quan lại phong kiến, truyền lại cho ông cách phân tích và suy ngẫm thời cuộc với một đầu óc tỉnh táo. Lớn lên khi học trung học và đại học ông tiếp xúc với những thầy cô phong cách cởi mở, dân chủ, tiên tiến. Khi ra làm việc ông lại được tiếp xúc với những trí thức ưu việt như Tạ Quang Bửu , Trần Đại NGhĩa. Khi làm nghiên cứu sinh ở Nga, ông lại tiếp xúc với những vị giáo sư giỏi về chuyên môn và rất sáng suốt về  ý thức hệ. Những bạn bè trong nước và ngoài nước của ông về sau phần lớn có dầu óc cởi mở, yêu chuộng cái mới. Trong lĩnh vực toán học, ít người biế được sự hy sinh to lơn và xót xa của ông khi từ bỏ lĩnh vực lí thuyết thuần túy ( siêu toán-métamathématique, toán kiến thiết-mathématique  constructive) để ngiên cứu và dần dắt giới trẻ làm công nghệ thông tin. Hãy hình dung nhà thơ Vũ Hoàng Chương từ bỏ thơ say để làm các vần thơ cảm tác trước các phương tiện văn minh như máy bay, tàu thủy.... Phan Đình Diệu ý thức được sự quan trọng của cách mạng 4.0 nên vui vẻ thi hành nhiệm vụ mới. Tiếc thay trên phương hướng này ông gặp không ít cản trổ vì tính hẹp hòi của những nhà lãnh đạo khoa học. Dân chúng biết đến ông và quí trọng ông vì những phát biểu công khai, mạnh mẽ trên các diễn đàn Quốc hội và MTTQ về xây dựng thể chế chính trị và đường lối chính sách. Cuối những năm 70, đầu những năm 80 thế kỉ trước, trong một kỳ họp quốc hội, ông đề nghị các nhà lãnh đạo cao tuổi nên nghỉ mà nhường chỗ cho thế hệ mới. Ông nói đại ý: Các đồng chí vĩ đại vì đã đem lại hòa bình thống nhất cho đất nước, nhưng các đồng chí sẽ vĩ đại hơn nếu về nghỉ nhường chỗ cho thế hễ sau. Ta nhớ lại môi trường tư tưởng của những năm đó khi câu "Đã thắng Mỹ thì không có gì mà ta không làm được" luôn được nhắc đi nhắc lại trong các hội nghị. Với hiểu biết về lí thuyết hệ thống,, về lô gich , về điều khiển học, Phan Đình Diệu cho rằng hai mặt trận quân sự và kinh tế xã hội không thể giống nhau để một người có thể điểu hành tốt cả hai.

Năm 1986, khi cả nước bước vào sự nghiệp đổi mới, Phan Đình Diệu cảnh báo rằng việc đổi mới về kinh tế bằng cách xây dựng kinh tế thị trường mà không đổi mới về thể chế thì chỉ tạo một nền kinh tế thiếu lành mạnh, bị các nhóm lợi ích lũng đoạn.
Rõ ràng tư tưởng đó đã được ứng nghiệm trong thực tiễn kinh tế xã hội mấy năm qua.

Phan Đình Diệu còn có nhiều suy nghĩ đóng góp cho các lĩnh vực khoa học, triết học, giáo  dục. Nguyên câu chuyện nhỏ về đời tư của ông đã thành một bài học  quí về cải cách hành chính. Có lần Đảng bộ chỗ ông làm việc đề nghị ông viết đơn xin vào Đảng. Ông vui vẻ trả lời: Nếu Đảng thấy tôi xứng đáng là đảng viên thì cứ việc tổ chức kết nạp tôi, việc gì phải bảo tôi làm đơn xin. Nền hành chính của ta, kể cả việc hành chính trong Đảng, xưa nay theo kiểu quan liêu bao cấp. Nhà nước không hề chủ động làm gì nếu không có đơn xin của người dân. Đáng ra những việc như cấp các danh hiệu Bà mẹ anh hùng, Huân chương các loại, Nhà giáo ưu tú...phải do các cấp chính quyền chủ động làm mà không bắt người dân phải viết đơn. Cụ thân sinh của tôi tham gia cách mạng từ năm 1927 rồi liên tiếp làm việc cho cách mạng cho đến khi về hưu. Có lẽ cụ không có ý kiến gì  nên cho đến khi mất cụ không được huân chương gì hết. Bà con làng xã thắc mắc quá trời nên gia đình mới chạy cho cụ được Huân chương Độc lập hạng ba. Nếu ai cũng nghĩ như Phan Đình Diệu thì xã hội sẽ tươi đẹp biết bao! Chính cái ý tưởng đó là phương châm cho việc cải cách hành chính, từ "cơ chế xin-cho" đến nền hành chính phục vụ nhân dân.
Chân thành chúc người bạn thân thiết yên nghỉ thanh thản ở chốn vĩnh hằng.